% SPDX-License-Identifier: CC-BY-4.0
% Copyright (C) 2026 the Langlandscape team
% Stage 03 section of the submission report (proof outlines condensed from mathematics/chapters/s03; statements verbatim).
% Note: the literature pointers (chapter macro lit{stem}{loc}) of the chapter text (also inside the verbatim boxes) were mechanically
% converted to \cite[loc]{KEY} with the keys of report.bib; nothing else in the boxes was changed except \ref -> main.pdf numbers.
\subsection{Stage 03: weighted graphs, the definite Ribet sequence and the Ihara lemma}\label{ss:stage03}

Stage 03 realises Ribet's integral comparison of character groups combinatorially. It attaches to a finite weighted graph
the lattice $X(\G)$ and the finite group $\Phicomb(\G)$, checks them against Mazur's component group at level~$1$
($\Phicomb\cong\Z/\mathrm{num}\frac{q-1}{12}$), proves that $\Phicomb$ is Eisenstein, and proves the definite Ribet sequence as the
homology sequence of a \emph{vertex collapse} of the bipartite Brandt graph of stage 02. The surjectivity in that sequence is the
definite Ihara lemma, and its cokernel with coefficients is computed, with Hecke action, as the Eisenstein module $V_{U^1_p}$.
Finally the Betti side for $\GL_2$ (amalgam, Mayer--Vietoris, $\Psi^*\Psi$) is recorded; the $\GL_2$ Ihara lemma in Betti
form stays a conjecture. The stage feeds stage 04 (ordinary Ribet sequence), stage 05 (identification with character groups) and stage 06
(the level-raising module). In main.pdf it occupies pp.~76--107 (Sections 8--12).

\subsubsection*{Key definitions}
\begin{itemize}
\item \emph{Weighted graph} (main.pdf, Def.~8.1, p.~76). A finite graph $\G$ (loops and multiple edges allowed) with oriented edges
  $o,t:E\to V$, chain complex $\Z[E]\xrightarrow{\partial}\Z[V]$, $\partial e=t(e)-o(e)$, and weights $w_e\in\Z_{\ge1}$. Put
  $X(\G):=H_1(\G)=\ker\partial$, $\langle e,e'\rangle_w=w_e\delta_{ee'}$, and
  \[\Phicomb(\G):=\coker\big(X(\G)\to\Hom(X(\G),\Z),\ x\mapsto\langle x,\cdot\rangle_w\big).\]
  Since $\langle,\rangle_w$ is positive definite, $|\Phicomb(\G)|=\det G$ for any Gram matrix $G$ of a basis of $X(\G)$, and
  $\Phicomb$ is independent of the orientation \stproved{} (basic facts after Def.~8.1). For a semistable curve with dual graph $\G$ and
  thicknesses $w_e$ this is the component group of the N\'eron model \cite[IX]{SGA7}, \cite[Prop.~2.2, Thm.~2.3, p.~436]{Ri90}
  \stcited{} (main.pdf, Remark~8.2); the tree route does \emph{not} use this dictionary.
\item \emph{DR graph} (main.pdf, Def.~8.6, p.~78). For $B=B_q$ (ramified at $\{q,\infty\}$) and $U=\Uone{Mp^s}$ or $\Uzero M$,
  $\GDR_q(U)$ is the two-vertex weighted graph with edge set $\XBq U$, every edge oriented $0\to1$, and
  $w_e=|\Gamma_e/(\Gamma_e\cap\{\pm1\})|$, $\Gamma_e=B^\times\cap xUx^{-1}$ for $e=[x]$. Then $X(\GDR_q(U))=\Z[\XBq U]^0$.
  For $s\ge1$ all $w_e=1$ and $\Phicomb\cong\Z/|\XBq{\Uone{Mp^s}}|$ \stproved{}; the interesting case is the maximal level $M=1$,
  $U=\hat\cO^\times$, where the weights are the half-orders of the unit groups of the maximal orders.
\item \emph{Bipartite Brandt graph and condition (iii)} (main.pdf, Def.~7.8, Prop.~7.9, stage 02). For $q\nmid D$ and a $q$-admissible
  level $U$ of $B=\Bd$ (Def.~7.4), let $U'=\Uzero qU$ be the Iwahori level at $q$, $\alpha=\iota_q^{-1}\mathrm{diag}(1,q)$, and
  $s_0[y]=[y]$, $s_1[y]=[y\alpha]$, $\XB{U'}\to\XB U$. $\BG(U,q)$ has vertex parts $V_0=V_1=\XB U$, edge set $\XB{U'}$, $o=s_0$, $t=s_1$.
  Condition (iii) is: \emph{$\BG(U,q)$ is connected}. It is a finite check at one level, and it is implied by
  $\HypSA\wedge\HypHS\wedge\mathrm{Hyp.BTtree}$. Under (iii), $\GU\backslash\cT_q\cong\BG(U,q)$ for the norm-one $S$-arithmetic group
  $\GU=\cO_U[\frac1q]^1$ (main.pdf, Prop.~7.10(a)).
\item \emph{Conventions (C1)--(C6)} (main.pdf, \S11.1, p.~91). (C1) edges are oriented from type $0$ to type $1$, so
  $\alpha:=(-s_{0*},s_{1*}):\Z[\XB{U'}]\to\Z[\XB U]^2$ and $\alpha^\vee(f_0,f_1)=s_1^*f_1-s_0^*f_0$ (Ribet's
  $\delta=(s_{0*},s_{1*})$ differs by $(x,y)\mapsto(-x,y)$). (C2) $w_x=|\Gamma_x/(\Gamma_x\cap\{\pm1\})|$, so the mass formula reads
  $\sum_xw_x^{-1}=\frac1{12}\prod_{\ell\mid D}(\ell-1)$. (C3) $s_{j*}$ on Brandt modules $\Brandt UV$ are the traces over $U/U'$, resp.\
  $U/\alpha^{-1}U'\alpha$. (C4) $(T_gf)(x)=\sum_if(xg_i)$ (cohomological). (C5) $\tau(x):=(t_{\Nrd x})_p\in\Zp^\times$, where
  $a=r_at_a$, $r_a\in\Q_{>0}$, $t_a\in\Zhat^\times$; then $\tau(\gamma xu)=\tau(x)\det(u_p)$, $\tau(x\alpha)=q^{-1}\tau(x)$.
  (C6) $U^1_p:=U_p\cap\SL_2(\Zp)$, and $[t]$ is the action of any $u\in U_p$ with $\det u=t$ on $V_{U^1_p}$ and on $V^{U^1_p}$.
\end{itemize}

\subsubsection*{Main results and proof outlines}

\paragraph{1. Two-vertex graphs and the Mazur check.} Everything at level $1$ reduces to the following linear algebra.

\begin{quoted}{Proposition}{8.3}{76}{prop:s03-phi}{two-vertex graphs; seed Prop.~phi}{chapters/s03/1-graphs.tex}\label{rpt:prop:s03-phi}
Let $\G$ have two vertices $v_0\ne v_1$ and $E\ne\emptyset$, and let every edge be oriented from $v_0$ to $v_1$. Then $X(\G)=\Z[E]^0$ and $\Phicomb(\G)=\Z^E/(W\Z[E]^0+\Z\mathbf1)$, where
$W=\mathrm{diag}(w_e)$. Moreover
\[|\Phicomb(\G)|=\Big(\prod_ew_e\Big)\sum_e\frac1{w_e}=\sum_e\prod_{e'\ne e}w_{e'}.\]
If all $w_e=1$, then $\Phicomb\cong\Z/|E|$. In general $\Phicomb$ is not cyclic: $w=(2,2,2)$ gives $\Z/2\times\Z/6$.
For an arbitrary orientation of the edges the order formula and the isomorphism class of $\Phicomb$ are unchanged (fact~(3)), but $X(\G)$ is
$\Z[E]^0$ only after the sign change of the reversed edges. For $E=\emptyset$ one has $\Phicomb=0$ of order $1$ while the sum is $0$, so the
formula fails; this case is excluded.
\stproved{} (proof below; finite linear algebra). Lean (S03-L3): \lean{Langlandscape.Graph.twoVertex\_X\_eq\_deg0},
\lean{Langlandscape.Graph.phiComb\_twoVertex\_equiv}, \lean{Langlandscape.Graph.card\_phiComb\_twoVertex},
\lean{Langlandscape.Graph.phiComb\_twoVertex\_unit} \stformal{}, for $E\ne\emptyset$ and all edges oriented $v_0\to v_1$ (the Lean hypotheses
\texttt{Nonempty G.E}, \texttt{hsrc}, \texttt{htgt}). The orientation-independence step (fact~(3)) and the example $w=(2,2,2)$ are not formalized.
\end{quoted}
\begin{outline}
\begin{enumerate}[label=(\arabic*)]
\item $\partial(\sum x_ee)=(\sum x_e)(v_1-v_0)$, so $X(\G)=\Z[E]^0$. \stproved
\item Restriction $\Z^E\to\Hom(\Z[E]^0,\Z)$ is onto with kernel $\Z\mathbf1$, and $x\mapsto\langle x,\cdot\rangle_w$ becomes $x\mapsto Wx$; this gives the
  presentation. \stproved
\item In the basis $b_i=e_i-e_n$ the Gram matrix is $D+w_n\mathbf u\mathbf u^{\mathsf T}$ with $D=\mathrm{diag}(w_1,\dots,w_{n-1})$; the matrix
  determinant lemma gives $\det G=\prod_iw_i\sum_iw_i^{-1}$. For $w=(2,2,2)$, $G=\left(\begin{smallmatrix}4&2\\2&4\end{smallmatrix}\right)$ has Smith form $\mathrm{diag}(2,6)$. \stproved
\end{enumerate}
More generally $|\Phicomb(\G)|=\sum_T\prod_{e\notin T}w_e$ over spanning trees of a connected $\G$ (weighted matrix-tree theorem, by Cauchy--Binet;
main.pdf, Remark~8.5, \stproved).
\end{outline}

\begin{quoted}{Lemma}{8.4}{77}{lem:s03-cyclic}{cyclicity}{chapters/s03/1-graphs.tex}\label{rpt:lem:s03-cyclic}
In the situation of Prop.~8.3, put $W:=\prod_ew_e$, $n:=|\Phicomb(\G)|$ and $\lambda(x):=\sum_e(W/w_e)x_e$. Then
$\lambda$ induces a homomorphism $\bar\lambda:\Phicomb(\G)\to\Z/n$. It is surjective, hence an isomorphism, if and only if
$\gcd_e(W/w_e)=1$. This holds in particular when the $w_e$ are pairwise coprime (any number of $1$'s allowed). \stproved{} (proof below;
the stage tex has \sttoprove). Lean (S03-L4): \lean{Langlandscape.Graph.lambdaBar},
\lean{Langlandscape.Graph.lambdaBar\_surjective\_iff}, \lean{Langlandscape.Graph.lambdaBar\_bijective\_of\_pairwise\_coprime} \stformal{}
(for two-vertex graphs with $E\ne\emptyset$, edges $v_0\to v_1$), with $\lambda$ itself as \lean{Langlandscape.Graph.lambdaMap} \stformal{} (a Lean definition on $\Z[E]$ with the
same formula).
\end{quoted}
\begin{outline}
$\lambda(\mathbf1)=W\sum_ew_e^{-1}=n$ and $\lambda(Wx)=W\sum_ex_e=0$ on $\Z[E]^0$, so $\lambda$ descends. Its image is $d\Z/n\Z$ with
$d=\gcd_e(W/w_e)$, of order $n/d$; a surjection between groups of order $n$ is bijective. If a prime $r$ divided $d$, it would divide some
$w_{e_0}$ and hence (pairwise coprimality) not $W/w_{e_0}$. \stproved
\end{outline}

\noindent\emph{The elliptic classes} (main.pdf, Lemma~8.7, p.~78, \stpmc). For $q\ge5$, $B=B_q$, $U=\hat\cO^\times$: every $w_x\in\{1,2,3\}$,
$n_2:=\#\{x:w_x=2\}=\frac12(1-(\frac{-4}q))$ and $n_3:=\#\{x:w_x=3\}=\frac12(1-(\frac{-3}q))$. Outline: by the list of finite unit groups of
quaternion orders \cite[Thm~11.5.14, p.~176]{Vo} \stcited, $\cO_x^\times$ is cyclic of order $2,4,6$ or binary dihedral of order $12$ or binary
tetrahedral of order $24$; the last two force $B\cong(-3,-1)_\Q$ resp.\ $(-1,-1)_\Q$ (explicit quaternion relations, Hilbert symbols), ramified at
$\{3,\infty\}$ resp.\ $\{2,\infty\}$, impossible for $q\ge5$ \stproved. The counts follow by subtracting the mass formula
$\sum_xw_x^{-1}=\frac{q-1}{12}$ (main.pdf, Thm.~2.1(a); \cite[Thm~25.1.1, p.~415]{Vo}) from the Eichler class number formula
$h=\frac{q-1}{12}+\frac{\epsilon_2}4+\frac{\epsilon_3}3$ (\cite[Thm~30.1.5, p.~543]{Vo}), both \stcited: $3n_2+4n_3\in\{0,3,4,7\}$ has a unique
solution. In Lean the statement is the registered hypothesis \pathtt{Hyp.EllipticCount}.

\begin{quoted}{Proposition}{8.8}{79}{prop:s03-mazur}{Mazur check; seed Example mazur}{chapters/s03/1-graphs.tex}\label{rpt:prop:s03-mazur}
For every prime $q$,
\[\Phicomb\big(\GDR_q(\hat\cO^\times)\big)\ \cong\ \Z\big/\mathrm{num}\tfrac{q-1}{12},\]
and this group is cyclic, generated via $\bar\lambda$ of Lemma~8.4. For $q\in\{2,3\}$ both sides are trivial.
\stpmc: Lemma~8.7, Prop.~8.3, Lemma~8.4. Lean (S03-L5, L6):
\lean{Langlandscape.Graph.mazurNum}, \lean{Langlandscape.Graph.mazur\_prod\_mul\_sum}, \lean{Langlandscape.Graph.mazur\_weights\_pairwise\_coprime}
\stformal{} (abstract arithmetic core: for weights in $\{1,2,3\}$ with the counts $n_2,n_3$ and $\sum w^{-1}=\frac{q-1}{12}$ one has
$\prod w\cdot\sum w^{-1}=\mathrm{num}\frac{q-1}{12}$, and the weights are pairwise coprime; $q\ge5$ only); the identification of the weights with the DR graph
(\lean{Langlandscape.Ribet.sum\_inv\_wx\_Ohat}, \lean{Langlandscape.Ribet.mazur\_check}, \stformal{} modulo the registered hypothesis \lean{Langlandscape.Hyp.EllipticCount} (and \pathtt{Hyp.EichlerMass})); the cases $q=2,3$ are \stformal{} as well, through the $q\le3$ clause of the same hypothesis \lean{Langlandscape.Hyp.EllipticCount} (which asserts $|\XB{\hat\cO^\times}|=1$ there; \lean{Langlandscape.Ribet.mazur\_check} covers every prime $q$), so the chapter's elementary derivation of that case is not formalized
\end{quoted}
\begin{outline}
\begin{enumerate}[label=(\arabic*)]
\item ($q\ge5$.) By Lemma~8.7, $\prod_ew_e=2^{n_2}3^{n_3}=12/g$ with $g=\gcd(q-1,12)$ (check $q\bmod4$ and $q\bmod3$). \stpmc
\item Prop.~8.3 and the mass formula: $|\Phicomb|=\prod w_e\cdot\sum w_e^{-1}=\frac{12}g\cdot\frac{q-1}{12}=\frac{q-1}g=\mathrm{num}\frac{q-1}{12}$. \stpmc{} (mass formula cited)
\item The weights lie in $\{1,2,3\}$ with $2$ and $3$ each occurring at most once, hence are pairwise coprime; Lemma~8.4 gives
  $\bar\lambda:\Phicomb\xrightarrow\sim\Z/\mathrm{num}\frac{q-1}{12}$. \stproved
\item ($q\in\{2,3\}$.) The mass is $\frac1{12}$ resp.\ $\frac16$, each $w_x\le12$ (resp.\ $\le6$ for $q=3$, since order $24$ forces ramification at $2$),
  so there is one class; one edge gives $\Phicomb=0$, and $\mathrm{num}\frac1{12}=\mathrm{num}\frac16=1$. \stpmc
\end{enumerate}
\emph{The picture.} At $M=1$ the DR graph is two vertices joined by $h=|\XBq{\hat\cO^\times}|$ edges, one for each class of maximal orders of
$B_q$, of thickness $w_e=|\cO_e^\times|/2$. All but at most two edges are thin ($w_e=1$); at most one has thickness $2$ (present iff
$q\equiv3\bmod4$) and at most one thickness $3$ (iff $q\equiv2\bmod3$). Since $\prod w\cdot\sum w^{-1}$ is the sum over edges of the product of the
other thicknesses, the order is $\mathrm{num}\frac{q-1}{12}$, and coprimality of the thicknesses makes $\Phicomb$ cyclic, generated by
$\bar\lambda^{-1}(1)$. Data (\stcomp, \res{s03/T1.json}; main.pdf, Example~8.10): $q=11$: weights $(2,3)$, $\Z/5$; $q=23$: $(1,2,3)$, $\Z/11$;
$q=37$: $(1,1,1)$, $\Z/3$; $q=101$: $(1^8,3)$, $\Z/25$; $q=131$: $(1^{10},2,3)$, $\Z/65$.
Mazur--Rapoport--Deligne's component group of $J_0(q)$ at $q$ is cyclic of order $\mathrm{num}\frac{q-1}{12}$, generated by $(0)-(\infty)$
\cite[App.\ Thm A.1(b), p.~173]{Mz} \stcited; the definite statement agrees in order and group structure. Identifying the two groups needs the
dictionary of Remark~8.2 (supersingular points as edges, $e(i)=\#\mathrm{Aut}(E)/2=w_e$; \cite[Prop.~3.1, 3.3, pp.~439--440]{Ri90},
\cite[VI Th.~6.9(iii)]{DR}, \stcited) and is \emph{not} part of the tree route (main.pdf, Remark~8.9).
\end{outline}

\paragraph{2. $\Phicomb$ is Eisenstein.} The Hecke operators prime to $qM$ act on $\Phicomb$ through their Eisenstein eigenvalues.

\begin{quoted}{Proposition}{9.5}{83}{prop:s03-phieis}{$\Phicomb$ is Eisenstein}{chapters/s03/2-eisenstein.tex}\label{rpt:prop:s03-phieis}
Let $q$ be any prime, $U=\Uzero M$, $B=B_q$, $\ell\nmid qM$, and let $T_\ell$ act on $\Z[\XBq U]$ by the Brandt matrix (homological
convention, $[x]\mapsto\sum_i[xg_i]$). Then $T_\ell$ is self-adjoint for $\langle,\rangle_w$, so it acts on $\Phicomb(\GDR_q(U))$, and
$T_\ell-(1+\ell)$ annihilates $\Phicomb(\GDR_q(U))$. \stproved{} (proof below, steps (1)--(4); valid for all primes $q$, including $q\in\{2,3\}$, and all
$M$ with $q\nmid M$ and $\ell\nmid qM$; it uses Lemma~9.2, Prop.~3.3(a),(b), Prop.~8.3, and the
finiteness of $\XBq U$, Prop.~1.6, which makes $\GDR_q(U)$ a finite graph). Lean: abstract core \stformal{}: \lean{Langlandscape.Graph.heckeQ} and \lean{Langlandscape.Graph.phiComb\_eisenstein} (steps (2)--(4), abstract: an integer matrix $B$ with row sums $\ell+1$,
$w_{e'}B(e,e')=w_eB(e',e)$ and $w_e\mid B(e,e')$ off the diagonal, acting on $\Phicomb$ of the two-vertex lattice $\Z[E]^0$ with weights $w_e>0$). Instantiation at $U=\hat\cO^\times$: \lean{Langlandscape.Ribet.phiComb\_eisenstein\_Ohat} \stformal{} modulo the registered hypothesis (step (1)) \lean{Langlandscape.Hyp.BrandtStabFree}.

Cross-references (not needed for the proof): [Ri90, Thm~3.12, p.~449; proof pp.~449--450, Remark~3.13] \stcited{}
\cite[Thm 3.12, p.~449]{Ri90}. Ribet's $\Phi=\coker(\iota)$ is our
$\Phicomb$, and his proof is the divisibility of the off-diagonal Brandt entries $B_\ell(i,j)$ by $e(i)=w_i$, which is step~(1) below, a purely quaternionic
fact. The geometric analogue for $\Phi_{N,p}$ is [Ed91, Th.~1, p.~163] \stcited{}
\cite[Edixhoven, Th\'eor\`eme~1, p.~163]{Ed91}, which assumes $p>3$.
\end{quoted}
\begin{outline}
Write $B(e,e')$ for the Brandt matrix, $T_\ell[e]=\sum_{e'}B(e,e')[e']$.
\begin{enumerate}[label=(\arabic*)]
\item \emph{$w_e\mid B(e,e')$ for $e\ne e'$} (main.pdf, Lemma~9.2). Neighbours of $x$ are the index-$\ell$ sublattices of $L_x=x_\ell\Z_\ell^2$
  (Lemma~9.1, \stproved), and $\Gamma_x/\{\pm1\}$ acts on them preserving classes. If $\gamma\ne\pm1$ fixes a neighbour $h$, then
  $\gamma$ has finite order, $\Nrd\gamma=1$, $\Trd\gamma\in\{-1,0,1\}$, so $A=\Z[\gamma]\in\{\Z[i],\Z[\omega]\}$ is a PID; $L_x$ is free of rank one over
  $\Z_\ell[\gamma]$, the fixed sublattice is $\pi L_x$ with $\pi\in A$, $\Nrd\pi=\ell$, and globally $hU=\pi xU$, so $[h]=[x]$. Hence the action on
  neighbours in other classes is free. \stproved{} (elementary; all $q$, all $\ell\nmid qM$, including $\ell\in\{2,3\}$; Lean: \pathtt{Hyp.BrandtStabFree}, discharged by \BSFdischargeName)
\item \emph{Symmetry} $w_{e'}B(e,e')=w_eB(e',e)$: from main.pdf, Prop.~3.3(a) ($T_g^*=T_{g^{-1}}$) and (b) ($S_\ell=1$ at level $\Uzero M$). This is self-adjointness of
  $T_\ell$ for $\langle,\rangle_w$; so $S(e,e'):=B(e,e')/w_e$ is a symmetric integer matrix off the diagonal. \stproved
\item On $\Z^E/\Z\mathbf1=\Hom(X,\Z)$, $((T_\ell^\vee-(\ell+1))\varphi)(e)=\sum_{e'\ne e}B(e,e')(\varphi(e')-\varphi(e))=w_ey_e$ with
  $y_e=\sum_{e'\ne e}S(e,e')(\varphi(e')-\varphi(e))$, and $\sum_ey_e=0$ by symmetry of $S$. \stproved
\item So $(T_\ell^\vee-(\ell+1))\Z^E\subseteq W\Z[E]^0$, which is the image of $X$; hence $T_\ell-(1+\ell)$ kills $\Phicomb$. \stproved
\end{enumerate}
Checked numerically for $\ell\le13$ in all rows of \res{s03/T1.json} \stcomp.
\end{outline}

\noindent\rptnum{9.7}\label{rpt:prop:s03-LRdef}\emph{The definite level-raising operator} (main.pdf, Prop.~9.7, p.~85, \stproved; Lean core and the instantiations at $\hat\cO^\times$ and
$\Uzero M$ \stformal). For $B=\Bd$, $q\nmid DM$, $U=\Uzero M$, the weighted adjoints $s_j^\dagger$ of $s_{j*}$ are integral,
$s_j^\dagger[x]=\sum_{u\in U/U'_j}[xua_j^{-1}]'$, and on $(\Z[\XB U]^0)^2$
\[\alpha\alpha^\dagger=\begin{pmatrix}q+1&-T_q\\-T_q&q+1\end{pmatrix},\qquad \det=(q+1)^2-T_q^2,\]
the definite analogue of Ribet's $\mu$ \cite[pp.~454--455]{Ri90}. Outline: $s_{j*}s_j^\dagger=[U:U'_j]=q+1$; $s_{1*}s_0^\dagger=T_\alpha=T_q$ since
$U\alpha U=\coprod_{u\in U/U'}u\alpha U$; $s_{0*}s_1^\dagger=T_{\alpha^{-1}}=T_q^\dagger$, which equals $T_q$ at level $\Uzero M$ because $z_q$ acts trivially.
Example (main.pdf, Example~9.8, \stcomp, \res{s03/T2.json}): $B_{11}$, $q=2$, $M=1$: $T_2=-2$ on $\Z[\XB U]^0$, $\det\alpha\alpha^\dagger=9-4=5$,
which is also $|\Phicomb|$ of the collapsed two-vertex graph with edge weights $(2,1,1)$.

\paragraph{3. The definite Ribet sequence as a vertex collapse.} The combinatorial core contains no quaternions.

\begin{quoted}{Lemma}{10.1}{87}{lem:s03-bip}{bipartite collapse}{chapters/s03/3-ribet.tex}\label{rpt:lem:s03-bip}
Let $\G$ be a finite \emph{connected} bipartite multigraph with parts $V_0,V_1$, edge set $E$, and maps $s_j:E\to V_j$. Put
$\partial=(-s_{0*},s_{1*}):\Z[E]\to\Z[V_0]\oplus\Z[V_1]$. Then:
\begin{enumerate}[label=(\alph*)]
\item $\ker\partial\subseteq\Z[E]^0$, and $\partial$ restricts to a surjection $\Z[E]^0\twoheadrightarrow\Z[V_0]^0\oplus\Z[V_1]^0$.
  Hence there is an exact sequence $0\to H_1(\G)\to\Z[E]^0\to\Z[V_0]^0\oplus\Z[V_1]^0\to0$, and $\rk H_1=|E|-|V_0|-|V_1|+1$.
\item For every abelian group $A$, the kernel of $\partial^\vee:A^{V_0}\oplus A^{V_1}\to A^E$, $(f_0,f_1)\mapsto s_1^*f_1-s_0^*f_0$, is
  the diagonal $\{(c\mathbf1,c\mathbf1):c\in A\}$.
\item If $E\ne\emptyset$ (equivalently, $V_0,V_1\ne\emptyset$), the sequence in (a) is the long exact homology sequence of the collapse
  $\G\to\G^{(2)}$ onto the graph with the same edges and two vertices. The connecting map is $\partial|_{\Z[E]^0}$. (If $E=\emptyset$, $\G$ is a
  single vertex and the sequence in (a) is $0$.)
\end{enumerate}
Parts (a) and (c) hold verbatim with $\Z$ replaced by any commutative ring $R$ ($R[E]$, $R[V_j]^0$, $H_1(\G,R)=\ker\partial$); the rank formula holds
for $R=\Z$ and for $\rk_R$ over any integral domain $R$. \stproved{} (proof below; the stage tex has \sttoprove; the seed's proof of Thm ribet(a)
is (c)). This is the Lean target: no quaternions are involved. Lean (S03-L2, any \texttt{CommRing R}): \lean{Langlandscape.Graph.bip\_ker\_le\_deg0}, \lean{Langlandscape.Graph.bip\_surj\_deg0}, \lean{Langlandscape.Graph.bip\_exact}, \lean{Langlandscape.Graph.bip\_rank} (over $\Z$), \lean{Langlandscape.Graph.bip\_ker\_dual} (any \texttt{AddCommGroup A}), \lean{Langlandscape.Graph.bip\_collapse\_connecting}, \lean{Langlandscape.Graph.bip\_collapse\_exact}, \lean{Langlandscape.Graph.bip\_collapse\_surjective} \stformal{}
\end{quoted}
\begin{outline}
\begin{enumerate}[label=(\arabic*)]
\item \emph{Image of $\partial$.} $\partial e=[t(e)]-[o(e)]$, so $\mathrm{aug}\circ\partial=0$; conversely each $[v]-[v_*]$ is a telescoping sum of $\pm\partial e$
  along a path, by connectedness. So $\partial(R[E])=\ker(\mathrm{aug})$. \stproved
\item (a) $s_{1*}$ sends basis vectors to basis vectors, so $\deg y=\deg s_{1*}y$; this gives $\ker\partial\subseteq R[E]^0$ and the surjection onto the
  degree-zero parts. The rank formula is additivity of rank on $0\to H_1\to\Z[E]\to\Z[V_0]\oplus\Z[V_1]\to\Z\to0$. \stproved
\item (b) $\partial^\vee(f_0,f_1)=0$ iff the function $F=f_0\sqcup f_1$ on $V_0\sqcup V_1$ is constant along edges, iff $F$ is constant (connectedness). \stproved
\item (c) The collapse $\pi:\G\to\G^{(2)}$ (identity on edges, $V_j\mapsto j$) is surjective on chains with kernel $K_\bullet$, $K_1=0$,
  $K_0=R[V_0]^0\oplus R[V_1]^0$. In the long exact sequence $H_1(\G^{(2)})=R[E]^0$, $H_0(K)=K_0$, $H_0(\G)\to H_0(\G^{(2)})$ is the identity of $R$, and the connecting
  map is $y\mapsto\partial y$. \stproved
\end{enumerate}
\end{outline}

\begin{quoted}{Theorem}{10.2}{89}{thm:s03-ribet}{definite Ribet sequence as a vertex collapse; seed Thm ribet}{chapters/s03/3-ribet.tex}\label{rpt:thm:s03-ribet}
Let $q\nmid D$, $U$ as in \texttt{if:s02-quot}, $\G:=\GU\quot\cT_q$, and $|\G|$ its underlying graph, with the orientation (C1).
\begin{enumerate}[label=(\alph*)]
\item Over $\Z$ there is an exact sequence, equivariant for all Hecke operators $T=[UgU]$ with $g$ supported at a finite set $\Sigma\not\ni q$ of
  places at which $U$ is factorisable (Prop.~7.10(c); this holds for all standard levels):
  \[0\to H_1(|\G|,\Z)\to\Z[\XB{U'}]^0\xrightarrow{\ \alpha=(-s_{0*},\,s_{1*})\ }\big(\Z[\XB U]^0\big)^{2}\to0.\]
\item Let $p\ge5$ and let $V$ be a finitely generated $\Zp$-module with continuous $U_p$-action. Then, with the traces of (C3),
  \[0\to H_1(\GU,V)\to\Brandt{U'}V\xrightarrow{\alpha}\Brandt UV^{2}\to H_0(\GU,V)\to0.\]
\item $\rk_\Z H_1(|\G|)=|\XB{U'}|-2|\XB U|+1$. At a level where all $\Gamma_x$ are trivial, this is $(q-1)|\XB U|+1$, as in
  \pathtt{if:s02-indef}. For $V=\Zp$, $p\ge5$: $H_1(\GU,\Zp)=H_1(|\G|,\Z)\otimes\Zp$.
\end{enumerate}
Part (a) holds over every commutative ring $R$ in place of $\Z$ (with $H_1(|\G|,R)$).
\emph{Status.} Given (iii): (a), (c) first two sentences \stproved{} (Lemma~10.1 and Prop.~7.10); (b) and the last
sentence of (c) \stpmc{} (Prop.~7.18(b), which rests on the cited tree property and Shapiro's lemma; Lean \pathtt{Hyp.SerreCoh}). Without
assuming (iii): \stpmc{} via $\HypSA\wedge\HypHS\wedge\mathrm{Hyp.BTtree}\Rightarrow$(iii) (Prop.~7.10). The definite form of Ribet's
surjectivity is [Ri90, Thm~3.15, p.~451]
\cite[Thm 3.15, p.~451]{Ri90} \stcited. Ribet proves it with Eichler's
theorems. [Ri90, Rem.~3.18, p.~452]
\cite[Rem.~3.18, p.~452]{Ri90} records the graph-connectivity
proof (Livn\'e, Edixhoven), which is ours. Lean (S03-L9, any \texttt{CommRing R}, under \pathtt{Tree.BGConnected}, or under $\HypSA\wedge\HypHS\wedge\mathrm{Hyp.BTtree}$ in the \texttt{\_of\_SA\_HS} versions): \lean{Langlandscape.Ribet.ribet\_exact}, \lean{Langlandscape.Ribet.ribet\_rank}, \lean{Langlandscape.Ribet.ribet\_hecke\_equivariant}, \lean{Langlandscape.Ribet.boundary\_BG\_eq}, \lean{Langlandscape.Ribet.boundary\_BG\_apply}, \lean{Langlandscape.Ribet.ribet\_exact\_of\_SA\_HS}, \lean{Langlandscape.Ribet.ribet\_rank\_of\_SA\_HS} \stformal{} for (a) and the rank in (c);
coefficient layer for (b), definitions and compatibilities only (S03-L14): \lean{Langlandscape.Ribet.CoeffData}, \lean{Langlandscape.Ribet.CoeffData.ofLocal},
\lean{Langlandscape.Ribet.alphaV}, \lean{Langlandscape.Ribet.alphaVee}, \lean{Langlandscape.Ribet.s0TrV\_comp\_s0PullV}, \lean{Langlandscape.Ribet.alphaV\_triv},
\lean{Langlandscape.Ribet.alphaVee\_triv} \stformal{}; the exactness statement (b) with coefficients $V$ itself is not formalized (\leanplanned)
\end{quoted}
\begin{outline}
\emph{The vertex-collapse picture.} The quotient $\GU\backslash\cT_q$ is the bipartite Brandt graph: one type-$0$ and one type-$1$ vertex for each
class of $\XB U$, one edge for each class of the Iwahori level $\XB{U'}$, joining $s_0[y]$ to $s_1[y]$. Collapsing all type-$0$ vertices to
one point and all type-$1$ vertices to another produces a two-vertex graph $\G^{(2)}$ with edge set $\XB{U'}$ and $H_1(\G^{(2)})=\Z[\XB{U'}]^0$; the
kernel of the collapse on chains is concentrated in degree $0$ and equals $(\Z[\XB U]^0)^2$. The homology sequence of the collapse is the Ribet
sequence: $H_1$ of the quotient graph (the tree model of the indefinite stage, main.pdf, Def.~7.21) is the kernel of the boundary map from the
``old'' edge lattice to two copies of the degree-zero vertex lattice, and the boundary map is onto.
\begin{enumerate}[label=(\arabic*)]
\item (a) Under (iii), main.pdf, Prop.~7.10(a) identifies $|\G|$ with $\BG(U,q)$, $o=s_0$, $t=s_1$; it is connected, and Lemma~10.1(a) gives the exact
  sequence with $\partial=\alpha$, over any $R$. \stproved{} given (iii)
\item \emph{Hecke equivariance.} $T=[UgU]$ acts at levels $U$ and $U'$ with the same coset representatives, preserves degree zero, and commutes with
  $s_{0*},s_{1*}$ (main.pdf, Prop.~7.10(c)); hence with $\alpha$. \stproved
\item (b) By main.pdf, Prop.~7.18(b), $H_*(\GU,V)$ is the homology of $\partial=(-s_{0*},s_{1*})$ on the homological Brandt modules
  \stpmc{} (tree property and Shapiro's lemma cited; Lean \pathtt{Hyp.SerreCoh}). For $p\ge5$ the norm maps are isomorphisms intertwining
  $s_{j*}$ with the traces (Prop.~7.14(c),(d)), which transports the sequence; the summand $V^-$ on which $-1$ acts by $-1$ contributes $0$ (Remark~7.15).
\item (c) Rank from Lemma~10.1(a); with trivial $\Gamma_x$, $|\XB{U'}|=(q+1)|\XB U|$ (Prop.~7.20(b)); the last claim is Prop.~7.22(b). \stproved{} / \stpmc
\end{enumerate}
$H_1(|\G|,\Z)$ differs from $H_1(\GU,\Z)$ by torsion at $2,3$ (main.pdf, Remark~10.4; e.g.\ $D=2$, $q=3$: $H_1(\GU,\Z)\cong\Z/3$, $H_1(|\G|,\Z)=0$).
Ribet's version for $\GL_2$ character groups is (a) together with two identifications this route does not prove here: \v{C}erednik--Drinfeld, and
$\HypX$ at a second prime (stage 05). Data (\stcomp, \res{s03/T2.json}; main.pdf, Example~10.3): $(D,q)=(2,7),(2,13),(11,2)$ give ranks $1,2,0$ in agreement with (c),
$\coker\alpha\cong\Z$ on the full modules, and $|\Phicomb|=6,21,5$ for the collapsed graph with the edge weights $(3,3)$, $(3,3,2)$, $(2,1,1)$.
\end{outline}

\paragraph{4. The definite Ihara lemma.} Surjectivity of $\alpha$ is the definite form of Ihara's lemma; with coefficients the cokernel is
$H_0$, which is Eisenstein.

\begin{quoted}{Corollary}{10.5}{90}{cor:s03-ihara}{definite Ihara lemma, (i)--(ii); seed Cor.~ihara}{chapters/s03/3-ribet.tex}\label{rpt:cor:s03-ihara}
(i) $\alpha$ in Thm~10.2(a) is surjective over $\Z$.
(ii) For every abelian group $A$ (in particular $A=\F_\ell$ for every prime $\ell$), the kernel of
$\alpha^\vee:(A^{\XB U})^2\to A^{\XB{U'}}$, $(f_0,f_1)\mapsto s_1^*f_1-s_0^*f_0$, is the line of constants $A\cdot(\mathbf1,\mathbf1)$.
(ii$'$) Let $p\ge5$ and let $W$ be a $\Zp$-module with continuous $U_p$-action. Then
\[\ker\big(\alpha^\vee:\Brandt UW^2\to\Brandt{U'}W\big)=\big\{(f_w,\,f_{[q]^{-1}w}):w\in W^{U^1_p}\big\},\qquad f_w(x):=[\tau(x)]^{-1}w,\]
and this is $H^0(\GU,W)=W^{U^1_p}$.
\stproved{} for (i), (ii) given (iii) (Lemma~10.1 and Prop.~7.10). (ii$'$): \stpmc{} (proof below, new: it checks
$f_w\in\Brandt UW$ by (C5) and $s_1^*f_{[q]^{-1}w}=s_0^*f_w$, and uses Lemma~7.17 (given (iii); the tree property is cited)
and the density of $\GU$ in $U^1_p$, i.e.\ Thm.~7.7; see (R4)). Here $U$ is assumed factorisable at $p$, $U=U_p\times U^{(p)}$ (true
for all standard levels), so that $\det:U_p\to\Zp^\times$ is onto and (C6) applies; ``continuous'' refers to a Hausdorff topology on $W$ for which
$W$ is a topological $\Zp$-module (e.g.\ $W$ finitely generated with the $p$-adic topology, or discrete). Lean (S03-L9): \lean{Langlandscape.Ribet.ihara\_surj}, \lean{Langlandscape.Ribet.ihara\_surj\_of\_SA\_HS} (i), \lean{Langlandscape.Ribet.ihara\_ker\_const}, \lean{Langlandscape.Ribet.ihara\_ker\_const\_of\_SA\_HS} (ii, any $A$) \stformal{}; (ii$'$) is not formalized.
\end{quoted}
\begin{outline}
\begin{enumerate}[label=(\arabic*)]
\item (i), (ii) are Lemma~10.1(a), (b) for the connected graph $\BG(U,q)$. \stproved{} given (iii)
\item (ii$'$), \emph{Step 1:} $f_w\in\Brandt UW$ for $w\in W^{U^1_p}$, by (C5) and since $u_p$ acts on $W^{U^1_p}$ as $[\det u_p]$. \emph{Step 2:}
  $s_1^*f_{[q]^{-1}w}(y)=[q^{-1}\tau(y)]^{-1}[q]^{-1}w=s_0^*f_w(y)$. \stproved
\item \emph{Step 3:} by main.pdf, Lemma~7.17 (given (iii); tree property cited) $\alpha^\vee$ is $F\mapsto F\circ\partial$ on $\GU$-equivariant cochains of the tree;
  as $0\to C_1\to C_0\to\Z\to0$ is exact, $\ker=W^{\GU}$ (constant cochains), which corresponds to $(f_w,f_{[q]^{-1}w})$ on tree representatives. \stpmc
\item \emph{Step 4:} $\GU$ is dense in $U^1_p$ (main.pdf, Thm.~7.7, full strong approximation) and stabilisers are closed, so $W^{\GU}=W^{U^1_p}$. \stpmc
\end{enumerate}
The dependence is explicit (main.pdf, table (R4), p.~91): (i), (ii) need only the finite connectivity check (iii); (ii$'$) and Thm.~11.3 need (iii) \emph{and} the density
of $\GU$ in $U^1_p$, which the per-level check does not provide.
\end{outline}

\begin{quoted}{Theorem}{11.3}{95}{thm:s03-iharaH0}{Ihara (iii): $H_0$ and its Hecke action}{chapters/s03/4-ihara.tex}\label{rpt:thm:s03-iharaH0}
Let $p\ge5$ and let $V$ be a finitely generated $\Zp$-module with continuous $U_p$-action. Use conventions (C1)--(C6).
\begin{enumerate}[label=(\alph*)]
\item The map
  \[\varepsilon:\Brandt UV^2\to V_{U^1_p},\qquad \varepsilon(f_0,f_1)=\sum_{[x]\in\XB U}w_x^{-1}\big([\tau(x)]\,f_0(x)+[q\,\tau(x)]\,f_1(x)\big)\]
  is well defined and surjective, and $\ker\varepsilon=\mathrm{im}\,\alpha$. Hence $\coker\alpha\cong H_0(\GU,V)=V_{U^1_p}$.
\item For $g$ supported at $\ell\nmid Dpq$, and $U$ factorisable at $\ell$, $U=U_\ell\times U^{(\ell)}$ (true for the standard levels,
  Def.~1.2), the induced action of $T_g$ on $\coker\alpha$ is
  \[T_g\ \longmapsto\ \deg(UgU)\cdot[\ell^{\,v_\ell(\Nrd g)}].\]
  In particular, $T_\ell$ acts as $(1+\ell)[\ell]$ and $S_\ell$ acts as $[\ell]^2$, the exponent being $+1$. If $U_p$ contains
  the scalars, then $[\ell]^2$ is the scalar $\ell$, which acts on $\Vk$ as $\ell^{k-2}$, in agreement with the classical
  normalisation $S_\ell=\ell^{k-2}$. The exponent $+1$ is for the cohomological operators $T_g$; for the transposed (homological) operators of the normative convention of stage 02, \pathtt{if:s02-indef}, it is $-1$, and there $S_\ell$ acts as $[\ell]^{-2}$ (Remark~11.4,
  exponent table, (R3)).
\item On the $\chi$-eigenpiece of $V_{U^1_p}\otimes\cOE$ for a character $\chi:\Zp^\times\to\cOE^\times$ (via $[\cdot]$), $T_\ell$ acts by
  $(1+\ell)\chi(\ell)$.
\end{enumerate}
\emph{Status.} \stpmc{} (proof below; the stage tex has \sttoprove{} with a sketch, which upgrades the seed's ``Hecke action \stsketch''). The
well-definedness of $\varepsilon$, $\varepsilon\circ\alpha=0$, surjectivity and the formula $\varepsilon\circ T_g=\deg(UgU)[\ell^m]\circ\varepsilon$ are \stproved{}
outright; $\ker\varepsilon\subseteq\mathrm{im}\,\alpha$ uses Cor.~10.5(ii$'$), hence (iii) and the density of $\GU$ in $U^1_p$ (full strong
approximation, (R4)); the identification with $H_0(\GU,V)$ is Prop.~7.18(b),(c). Lean (S03-L11): the definitions \lean{Langlandscape.Ribet.tau}, \lean{Langlandscape.Ribet.diamond}, \lean{Langlandscape.Ribet.epsilon} and the predicate
\lean{Langlandscape.Ribet.EpsWellDefined} \stformal{} (definitions only); clause (a) is the \emph{registered hypothesis} \lean{Langlandscape.Hyp.IharaH0} over
\pathtt{Ribet.alphaV} (an assumption, not a formalized theorem); clause (b) is
maths only here (the Hecke operators with coefficients, S04-L15a, now exist in Lean, but clause (b) itself is neither stated nor proved), \leanplanned
\end{quoted}
\begin{outline}
Two auxiliary lemmas, both \stproved: the \emph{unfolding identity} (main.pdf, Lemma~11.1) $\sum_{[x]\in\XB U}w_x^{-1}\sum_{u\in U/U''}F(xu)=\sum_{[y]\in\XB{U''}}w_y^{-1}F(y)$
for $-1\in U''\iff-1\in U$ (orbit counting, the factor-$2$ hypothesis of (C3) holds for $U'$, $U'_\alpha$, $U\cap gUg^{-1}$); and \emph{adjoints with coefficients}
(Lemma~11.2): for $\langle f,\varphi\rangle_U=\sum_xw_x^{-1}\langle f(x),\varphi(x)\rangle$, $s_{j*}$ is adjoint to $s_j^*$, $T_g^\dagger=T_{g^{-1}}$, the
pairing with $V^*=\Hom(V,\Qp/\Zp)$ is perfect, and $\deg(Ug^{-1}U)=\deg(UgU)$ (unimodularity).
\begin{enumerate}[label=(\arabic*)]
\item \emph{(a1) Well defined:} for $x'=\gamma xu$, $\tau(x')=\tau(x)\det u_p$ and $f_j(x')=u_p^{-1}f_j(x)$, and $u_p^{-1}$ acts on $V_{U^1_p}$ as $[\det u_p]^{-1}$. \stproved
\item \emph{(a2) $\varepsilon\circ\alpha=0$:} with $F(y)=[\tau(y)]\overline{g(y)}$, unfolding turns both components of $\varepsilon(-s_{0*}g,s_{1*}g)$ into $\sum_{[y]\in\XB{U'}}w_y^{-1}F(y)$
  (for $s_1$ using $\tau(xu\alpha^{-1})=q\tau(x)$, which is the reason for the factor $[q]$). \stproved
\item \emph{(a3) Onto:} the function supported on $[x]$ with value $N_xv$ maps to $[\tau(x)]\bar v$, since $H_x\subseteq U^1_p$ acts trivially on coinvariants. \stproved
\item \emph{(a4) $\ker\varepsilon=\mathrm{im}\,\alpha$:} by Pontryagin duality and Lemma~11.2, $(\mathrm{im}\,\alpha)^\perp=\ker(\alpha^\vee)$ on $\Brandt U{V^*}^2$, which
  by Cor.~10.5(ii$'$) with $W=V^*$ is $\{(f_\varphi,f_{[q]^{-1}\varphi})\}$; the pairing with these is $\varphi(\varepsilon(f_0,f_1))$. Hence
  $\ker\varepsilon=(\mathrm{im}\,\alpha)^{\perp\perp}=\mathrm{im}\,\alpha$. \stpmc{} (via (ii$'$): density of $\GU$ in $U^1_p$)
\item \emph{Identification} $\coker\alpha\cong H_0(\GU,V)=V_{U^1_p}$: Thm.~10.2(b) and main.pdf, Prop.~7.18(c). \stpmc
\item \emph{(b)} With $m=v_\ell(\Nrd g)$, $\tau(yg)=\ell^{-m}\tau(y)$; unfolding for $W=U\cap gUg^{-1}$ and $W'=g^{-1}Wg$ gives
  $\varepsilon(T_gf_0,T_gf_1)=[U:W'][\ell^m]\varepsilon(f_0,f_1)$ and $[U:W']=\deg(UgU)$. For $T_\ell$: $(1+\ell)[\ell]$; for $S_\ell$ ($g=z_\ell$, one coset): $[\ell]^2$.
  A duality cross-check on $H^0$ gives the same exponent. \stproved{} (as an identity for $\varepsilon$; via $\bar\varepsilon$ it inherits (a4))
\item \emph{(c)} is immediate from (b). \stproved
\end{enumerate}
\emph{Exponent.} For the homological (transposed) operators of stage 02 the exponent is $-1$: $T_\ell\mapsto(1+\ell)[\ell]^{-1}$, $S_\ell\mapsto[\ell]^{-2}$, so the
normative $T_\ell,S_\ell$ are the classical $S_\ell^{-1}T_\ell,S_\ell^{-1}$ (main.pdf, Remark~11.4 and the exponent table, \stproved). HACS test T6 ($B_3$, tree at
$q=2$, $p\in\{5,7\}$, $\Uone p$-level) distinguishes the two exponents and finds $+1$ (cohomological) and $-1$ (homological) as predicted \stcomp{},
\res{s03/T6.json}. At maximal level $U_p=\GL_2(\Zp)$ and $2<k\le p+1$, $H_0(\GU,\Vk)=0$ (irreducibility of $\Sym^{k-2}\F_p^2$ and Nakayama); at
$\Uone{p^s}$-level, $H_0\ne0$ for every $k$ (main.pdf, Remark~11.6(a),(b), \stpmc{} given Prop.~7.18(c)); the latter is the finite-level Eisenstein cokernel that stage 04 kills
with the ordinary projector.
\end{outline}

\begin{quoted}{Corollary}{11.5}{98}{cor:s03-eis}{Eisenstein support}{chapters/s03/4-ihara.tex}\label{rpt:cor:s03-eis}
Call a maximal ideal $\m$ of the Hecke algebra $\T$ (generated over $\Zp$ by the $T_\ell,S_\ell$ with $\ell\nmid NDpq$, where $N$ is the level of $U$ away from $Dpq$) \emph{Eisenstein}
if there is $a\in\Z/(p-1)$ such that $T_\ell-\ell^a(1+\ell)\in\m$ and $S_\ell-\ell^{2a}\in\m$ for all such $\ell$. This implies
Eisenstein in the sense of [DT, p.~439--440]. Then every $\m$ in the support of $H_0(\GU,V)$ is Eisenstein, and
$H_0(\GU,V)_\m=0$ for non-Eisenstein $\m$. \stpmc{} (proof below, from Thm~11.3). The classical analogue is [DT,
Lemma~2 and Cor.~1, pp.~445--446] \stcited{} \cite[Lemma 2, p.~445, and Cor.~1, p.~446]{DT} (Lemma~2 is the Ihara-type lemma, Cor.~1 the Eisenstein-support statement).
Lean (planned, optional, S03-L11): \pathtt{Ribet.eis\_support\_of\_IharaH0}. \leanplanned
\end{quoted}
\begin{outline}
By Thm.~11.3(b), $\T$ acts on $H_0\cong V_{U^1_p}$ through $\T\to A$, $T_\ell\mapsto(1+\ell)[\ell]$, $S_\ell\mapsto[\ell]^2$, where $A$ is the (finite, commutative) image of
$\Zp[\Zp^\times]$. For $\m$ in the support choose a maximal ideal $\mathfrak n$ of $A$ above it (lying over); $t\mapsto[t]\bmod\mathfrak n$ is a continuous character to a finite
field of characteristic $p$, trivial on the pro-$p$ group $1+p\Zp$, hence $t\mapsto\bar t^a$. So $T_\ell\equiv(1+\ell)\ell^a$, $S_\ell\equiv\ell^{2a}$ mod $\m$. \stpmc{}
(inherits the status of Thm.~11.3)
\end{outline}

\paragraph{5. The Betti side for $\GL_2$.} Let $(M,q)=1$, $\Gamma=\Gamma_0(M)$, $G=\Gamma_0(M)[\frac1q]$, $\alpha'=\mathrm{diag}(q,1)$,
$\Gamma'=\alpha'^{-1}\Gamma\alpha'$, so $\Gamma\cap\Gamma'=\Gamma_0(Mq)$; $l$ prime with $l>k$, $l\ne q$; $k\ge4$ even; $V=\Sym^{k-2}(\Z_{(l)}^2)$ (standard column
representation).
\begin{itemize}
\item \emph{Amalgam} (main.pdf, Prop.~12.2, \stpmc): $G=\Gamma*_{\Gamma_0(Mq)}\Gamma'$, by \cite[II.1.4 Thm~3, p.~79]{Se} \stcited{} applied to the action on $\cT_q$;
  Serre's hypothesis (D), density of $G$ in $\SL_2(\Qq)$, is proved elementarily from the unipotent subgroups (Lemma~12.1(d), \stproved).
\item \emph{Mayer--Vietoris} (main.pdf, Prop.~12.3, \stpmc; \cite[III.6 Exercise, p.~74]{Br} \stcited): $0\to H^1(G,V)\to H^1(\Gamma,V)^2\xrightarrow\Psi H^1(\Gamma_0(Mq),V)\to H^2(G,V)\to0$
  with $\Psi(x,y)=\mathrm{res}\,x-\mathrm{res}(cy)$. The end terms vanish because $\Sym^{k-2}\F_l^2$ has no invariants under $\Gamma_0(Mq)$ (also if $l\mid M$) and
  $\Gamma$ is virtually free with torsion prime to $l$; $H^1(\Gamma,V)$ is torsion-free. By Lemma~12.8, $H^1(G,V)=0$: the kernel of $\Psi$ is killed by
  $\det(\Psi^*\Psi)$, and no Hecke eigenvalue $\lambda$ of $T_q$ on $H^1(\Gamma,V\otimes\C)$ (cusp forms by Ramanujan--Petersson \cite[Thm.~8.2]{De74}, Eisenstein series
  by their explicit eigenvalues, both \stcited) satisfies $\lambda^2=q^{k-2}(q+1)^2$. \stpmc
\end{itemize}

\begin{quoted}{Proposition}{12.6}{104}{prop:s03-PsiPsi}{$\Psi^*\Psi$}{chapters/s03/5-amalgam.tex}\label{rpt:prop:s03-PsiPsi}
Put $\Psi^*(z):=\big(\mathrm{cor}^\Gamma_{\Gamma_0(Mq)}z,\ -c^{-1}\mathrm{cor}^{\Gamma'}_{\Gamma_0(Mq)}z\big)$. Then
\[\Psi^*\Psi=\begin{pmatrix}q+1&-T'_q\\-T''_q&q+1\end{pmatrix}\ \text{ on }H^1(\Gamma,V)^2,\qquad
T'_q=\mathrm{cor}\circ\mathrm{res}\circ c,\quad T''_q=c^{-1}\circ\mathrm{cor}'\circ\mathrm{res}.\]
Here $T'_q=T_{\alpha'^{-1}}$ is the operator of the double coset $\Gamma\alpha'^{-1}\Gamma=q^{-1}\Gamma\alpha'^\iota\Gamma$ (acting through its elements), and
$T''_q=T_{\alpha'}$ is that of $\Gamma\alpha'\Gamma$. For $V=\Sym^{k-2}$ of the standard (column) representation and $T_q:=T_{\alpha'}$, the classical
operator (Lemma~12.5),
\[T'_q=q^{2-k}T_q,\qquad T''_q=T_q,\qquad \Psi^*\Psi=\begin{pmatrix}q+1&-q^{2-k}T_q\\-T_q&q+1\end{pmatrix},\]
\[\det(\Psi^*\Psi)=(q+1)^2-q^{2-k}T_q^2=-q^{2-k}\big(T_q^2-q^{k-2}(q+1)^2\big),\]
which is Ribet's level-raising condition up to a unit. \stpmc{} (formal properties of $\mathrm{res},\mathrm{cor},c_g$ cited from Brown, Lemma~12.4;
the identification of $T_q$ with the classical operator is Lemma~12.5(d); the stage tex has \sttoprove). The determinant is \stcomp{}
(see after the table below).
\end{quoted}
\begin{outline}
\begin{enumerate}[label=(\arabic*)]
\item Diagonal: $\mathrm{cor}\circ\mathrm{res}=[\Gamma:\Gamma_0(Mq)]=q+1$ (Lemma~12.4(c), 12.1(b)). Off-diagonal: by definition. \stpmc{} (Brown cited)
\item Compatibility of $c_g$ with $\mathrm{res}$ and $\mathrm{cor}$ identifies $T'_q=T_{\alpha'^{-1}}$ and $T''_q=T_{\alpha'}$. \stpmc
\item $\alpha'^{-1}=q^{-1}\alpha'^\iota$, scalars act by $\lambda^{k-2}$, and $\Gamma\alpha'^\iota\Gamma=\Gamma\alpha'\Gamma$ (explicit matrices, Lemma~12.1(c)), so $T'_q=q^{2-k}T_q$. \stproved
\item $T_q=T_{\alpha'}$ is the classical Hecke operator: on $H^1_P$ via the Eichler--Shimura isomorphism \cite[\S6.3, Thm.~1, p.~177]{Hi93} \stcited, on the Eisenstein part by
  Lemma~12.5(e) \stpmc; checked on all of $H^1$ in 14 cases (HACS, \res{s03/C3\_rerun.json}, \stcomp).
\item The determinant is computed in the commutative ring $\Z_{(l)}[T_q]$; $\det$ and $\coker$ are independent of the normalisation (conjugation by $\mathrm{diag}(1,q^{2-k})$). \stproved
\end{enumerate}
$\det(\Psi^*\Psi)\ne0$ on each $T_q$-eigenvalue of $S_k(\Gamma_0(M))$ in all 121 cases of \res{s03/betti\_det.json} \stcomp. Which off-diagonal slot carries $q^{2-k}$ is
not tested numerically; it rests on the proof. The \emph{Betti level-raising module} $\mathrm{LR}_q(M,k)_\m:=\coker(\Psi^*\Psi)_\m$ (main.pdf, Def.~12.9) is finite whenever
$T_q^2-q^{k-2}(q+1)^2$ is a non-zero-divisor, which holds under the hypotheses of Lemma~12.8 \stproved; stage 06 uses it as the Betti incarnation of the
level-raising ideal. The classical source is \cite[\S2--3]{Ri84} \stcited.
\end{outline}

\begin{quoted}{Conjecture}{12.10}{107}{conj:s03-iharaGL2}{Ihara lemma for $\GL_2$, Betti form}{chapters/s03/5-amalgam.tex}\label{rpt:conj:s03-iharaGL2}
Let $\m$ be a maximal ideal of $\T^{(Mql)}$ (Lemma~12.7); call $\m$ \emph{Eisenstein} if there is $d\ge1$ such that $T_p-2\in\m$ for all but finitely
many primes $p\equiv1\bmod d$. This is the notion of [DT, p.~439--440]
\cite[p.~439--440]{DT}; DT's second condition $S_p-1\in\m$ is automatic for $p\equiv1\bmod dl$, since the
scalar $p$ acts as $p^{k-2}$ (Lemma~12.4(b)) and $l\in\m$. Compare Cor.~11.5 on the quaternionic side.
For non-Eisenstein $\m$, $(\coker\Psi)_{\m,\mathrm{tors}}=0$. Equivalently, $\Psi\otimes\F_l$ is injective after localisation at
$\m$, because the groups $H^1(\cdot,V)$ are torsion-free (so the Tor sequence of $\otimes\F_l$ has no extra term; $H^2(\Gamma,V/l)$ plays no role) and $\Psi$ is injective (Lemma~12.8). \stexp{} (consistent with C3: all torsion found is
Eisenstein, at $l=k+1$). The classical route uses Ihara's lemma [Ih75] (not in \pathtt{literature/}; \texttt{ihara\_1966} there is a different paper) or the congruence subgroup
property of $\SL_2(\Z[\frac1q])$ (\pathtt{literature/serre\_1970\_*}, scan, theorem not checked). The indefinite quaternionic analogue
is [MS19, Thm~1.1]
\cite[Thm 1.1, p.~2]{MS19} \stcited.
The ``equivalently'' needs $\ker\Psi=H^1(G,V)=0$, which is Lemma~12.8. Writing $0\to K\to A\xrightarrow\Psi B\to C\to0$ with
$A=H^1(\Gamma,V)^2$, $B=H^1(\Gamma_0(Mq),V)$ free over $\Z_{(l)}$ (Prop.~12.3), one has $\ker(\Psi\otimes\F_l)\supseteq K/l$ and
$\ker(\Psi\otimes\F_l)/(K/l)\cong C[l]$; so, as $K=0$, $\Psi\otimes\F_l$ is injective after localisation at $\m$ iff $C[l]_\m=0$ iff $C_{\m,\mathrm{tors}}=0$
(all maps are $\T^{(Mql)}$-linear by Lemma~12.7).
\end{quoted}
\begin{outline}
No proof is claimed; the statement is \stexp. What is proved is the equivalence of the two formulations (from $\ker\Psi=0$, Lemma~12.8, \stpmc, and
the freeness of $H^1(\Gamma_0(Mq),V)$) and the Hecke linearity of $\Psi$ for $p\nmid Mql$ (Lemma~12.7, \stpmc). The quaternionic counterpart, the
definite Ihara lemma (Cor.~10.5, Thm.~11.3, Cor.~11.5), \emph{is} proved on the tree route; the $\GL_2$ statement would follow classically from
\cite{Ih75} or the congruence subgroup property of $\SL_2(\Z[\frac1q])$, neither of which is checked here.
\end{outline}

\subsubsection*{What this stage does not prove}
\begin{itemize}
\item The $\GL_2$ Ihara lemma in Betti form (Conj.~12.10) is \stexp{}; only the quaternionic (definite) Ihara lemma is proved.
\item The identification of $\Phicomb(\GDR_q(\hat\cO^\times))$ with the component group of $J_0(q)$ (Remark~8.9), and of $H_1(\BG(U,q))$ with a character group of a
  Shimura curve or $J_0$ (\v{C}erednik--Drinfeld, $\HypX$; stage 05), are \stcited{} resp.\ outside this stage; the tree route does not use them.
\item The inputs not proved here: the mass and Eichler class number formulas and the list of unit groups (Lemma~8.7, \stcited); the tree property and
  Shapiro's lemma behind Prop.~7.18(b) (Thm.~10.2(b)); full strong approximation, i.e.\ density of $\GU$ in $U^1_p$ (Cor.~10.5(ii$'$), Thm.~11.3, Cor.~11.5);
  Serre's amalgam theorem, Brown's Mayer--Vietoris, Hida's Eichler--Shimura comparison and Deligne's Ramanujan--Petersson bound (Section 12); all \stcited.
\item Lean: Thm.~10.2(b) with coefficients, Cor.~10.5(ii$'$), Thm.~11.3(b) and Cor.~11.5 are not formalized; Thm.~11.3(a) is the registered hypothesis
  \texttt{Hyp.IharaH0}; Lemma~8.7 is \pathtt{Hyp.EllipticCount} and step (1) of Prop.~9.5 is \pathtt{Hyp.BrandtStabFree} (proved or cited on paper; the former is open in Lean, the latter discharged by \BSFdischargeName).
  None of Section 12 is formalized.
\item The off-diagonal slot of $q^{2-k}$ in $\Psi^*\Psi$ is not independently confirmed numerically (PARI/HACS test only $\det$).
\end{itemize}
Secondary results omitted above, with statuses: Lemma~9.1 (neighbours) \stproved; Remarks~9.3, 9.4 (decomposition of $\ell$; ideal language) \stproved{} resp.\
\stpmc; Lemma~12.4 \stpmc; Lemma~12.5 (a)--(c) \stproved, (d), (e) \stpmc; Lemma~12.7 \stpmc; Remark~11.6 (a), (b) \stpmc, (c) \stcomp.
